% !TeX root = ../papers/nc-p3.tex
\section{二項中心の二次錐と第二次数の強制}
\label{sec:binary-centers-quadric-cones}
\begin{theorem}\label{thm:binary-centers-quadric-cones}
曲面射影のない標準正則窓で二直線列がないとする。中央支持2の順中心 $p$ があれば、
その一つ前の位相に順中心はなく、その位相には $p$ しかない。
$p$ の第二曲線は必ず滑らかな平面二次曲線であり、最初の直線上のただ一つの順中心は
中央支持3である。第三曲線は平面二次またはねじれ三次で、第二曲線は
$p$ の中央支持の消去直線に長さ1で会う。逆向きにも成り立つ。
四位相を合わせても中央支持2の順中心は高々一つである。
これらは純関係や反対階数への仮定を要せず、特に反対階数10に適用できる。
\end{theorem}
\begin{proof}[錐と三次曲線の後の下界]
$p$ を位置0とする。定理\ref{thm:branch-binary-center-pairing}の対となる逆中心 $z$ は
位置3にあり、$L=L_p$ は位置1の後続直線、$N$ は位置2の中央支持の消去直線である。
$N$ の各点の後続は $z$ 一つである。$p$ の一意な前点を $p^-$ とする。
二項対称二乗の表示で、中央の二座標を $a,b$、位置3の末尾三座標を $c_0,c_1,c_2$ とすれば、
$R_2$ の二関係は $a c_0+b c_1, a c_1+b c_2$ である。従って位置2,3の点対で
位置2が $N$ の外、位置3が $p^-$ の外なら
\[
             [c_0:c_1:c_2]=[b^2:-ab:a^2].
\]
位置3の動く成分は全て頂点 $p^-$ のこの二次錐上にある。
実際前点が $N$ にあれば後続は $z$ 一点にすぎない。逆向きの対に同じことを使えば
位置0の動く成分も $z$ の一意な次点を頂点とする二次錐にある。
ここで用いる成分は全て列の成分で、前点と次点の存在が保証されている。
一つの位相で同じ二次錐上にある列の異なる成分の次数総和は8以下である。
射影評価の四次小行列式の少なくとも一つはこの錐上で非零であり、
その零点の因子が全成分を含むからである。

第二曲線 $x_2$ がねじれ三次曲線の場合を考える。その $N$ からの射影の次数を
$k=3-n$、$n=\operatorname{length}x_2^{-1}(N)$ とすれば $k>0$ で、
$x_3$ の $p^-$ からの射影は上述の二次写像との合成になる。
第三曲線は動き、その次数は $d_3=2k+m$、$m=\operatorname{length}x_3^{-1}(p^-)\le1$ である。
最後の不等式は $p^-\to p$ の固定帰路を付ける前節の閉経路の証明による。
$d_3\le3$ なら $k=1$、$n=2$ である。$x_3$ が二次なら滑らか、三次なら非平面的である
（二次錐の平面切断には平面三次曲線がない）。従って後者はねじれ三次で、どちらも閉埋込みである。
一方 $x_2^{-1}(N)$ 上で $x_3=z$ がスキームとして成立する。
実際 $N$ 上の順評価は階数3で、$z$ は全 $N$ で解である。局所の単元な三次小行列で解けば
任意の微小部分上でも核直線はこの定点となる。長さ2が一点に写るのは閉埋込みに反する。
従って $d_3\ge4$ であり、特に二項中心側からは $(1,3,1)$ の列は出ない。
\end{proof}
\begin{proof}[一つ前の位相]
$R_1$ の三関係は二項対称二乗なので、位置1が $L$ の外の点対は位置2が $N$ に入る。
従って位置1の順中心は全て $L$ 上である。外にあればその後続直線全体が $N$ で、
ただ一つの動く位置の後に $z$ へ至ってしまう。位置0に他の順中心があればその後続直線は $L$ と交わらず、
第二曲線が $N$ となり二直線列の排除に反する。
位置3に順中心があると仮定する。その第二曲線は非直線（第一直線の後も直線なら
次数式で二直線列となる）なので $L$ ではない。その第三曲線は $N$ であり、
次は $z$ なので次数列は $(1,d,1)$、$d=2,3$ である。
逆に $z$ から始めて前段を使うと $d=3$ は不可能である。

$d=2$ なら次数式により第二曲線での順側欠損は $3\cdot2-3-1=2$。
その上のある順中心 $r$ の後続直線は $N$ に会う。一方 $r\in L$ だから $w(r,z)=0$、
従ってその中央支持は2または3である。2の場合は、$p$ の中央支持の各文字に終わる純関係と、
$r$ の順支持の各文字に始まる純関係により、両二空間は交われない。
従ってその消去直線も交わらない。3の場合、$L_r$ の一般中央支持は4で、
順中心でない全点でも4である。実際そこでは定理\ref{thm:finite-branch-corridors}の
三列が独立で中央支持が落ちない。二項中心の対の定理から中央支持3の最初の直線は二項中心から出るものに限る。
$N$ 上の全点は順中心でなく、反対位置の $p$ に消され中央支持が高々3なので、やはり交われない。
これで位置3には順中心がない。
\end{proof}
\begin{proof}[二次錐の次数と第二曲線]
再び $d_2=3$ と仮定する。位置1の全順中心は $L$ 上にあり、その上の欠損 $f_1=0$ なので存在しない。
位置2の順中心の後続直線は全て位置3の二次錐上に入る。全母線は頂点を通り、
異なる中心の後続直線は交わらないため、この位相の順中心は高々一つである。
一つもなければ $f_2=0$ より $d_3=6$ となり、$f_3=0$ から
$d_4=3d_3-3d_2+1+b_2\ge10$ となる。第三曲線は非直線なので第四曲線も動き、
位置0の二次錐の上界8に反する。

位置2の順中心を $a$ とすれば、その後続直線に順中心はないのでその第二曲線は
位置0を張る次数3の曲線である。$p$ の第四曲線はこれと異なる
（異なる列が同じ動く成分を使えない）。従って $d_4\le8-3=5$。
$d_4=3d_3-8+b_2$ と $d_3\ge4$ から $d_3=4$、$d_4\ge4$ である。
$d_3=2k+m$ より $m=0$、$k=2$、$n=1$。
第三曲線は $p^-$ を避けるため第四曲線には $p$ がなく、位置0の順中心の一意性から $f_4=0$。
一方 $n=1$ から第三曲線は逆中心 $z$ を通り $b_3\ge1$ である。
従って $d_5=3d_4-3d_3+d_2+b_3-f_4\ge4$ となる。
この第五曲線は位置1の $L$ と異なり、非直線なので第六曲線も動いて $N$ に一致する。
次は $z$ であり、逆向きに見ると第五曲線は第二曲線となる。しかし第二曲線の次数は3以下である。
これで $d_2=3$ を排除した。

従って $d_2=2$ で、$f_1=1$ より第一直線に順中心が一つある。
それは $z$ に消され中央支持3以下であり、中央支持2なら前段の一位相前の禁止に反する。
二次曲線から $N$ への交叉長を再び $n$ とする。
$N$ からの射影が一定ならこの曲線の平面が $N$ を含み、$n=2$ である。
この場合第三曲線の射影も一定なので第三曲線は直線（または定点だが二直線列の排除から動く）であり、
先の長さ2の矛盾が成立する。従って射影次数 $k=2-n$ は正である。
$d_3=3-f_2\le3$ と $d_3=2k+m$ より $k=1,n=1$ であり、先と同じ錐の平面切断によって
第三曲線は平面二次またはねじれ三次である。
もう一つ二項順中心があれば相対する位置2に限られるが、その前位相に上記の三項順中心があるため矛盾する。
\end{proof}

\begin{corollary}\label{cor:no-binary-centers-without-two-lines}
曲面射影がなく、どの列も二直線列でない標準正則窓では、順中心・逆中心の中央支持は
全て3または4である。全ての動く成分の一般点の中央支持は4である。
特に反対階数10の曲線分岐に中央支持2の候補はない。
\end{corollary}
\begin{proof}
上の二項順中心 $p$ とその対の逆中心 $z$ があるとする。定理とその逆により
位置3に順中心はなく、位置0に逆中心はない。
位置2に順中心 $a$ があればその最初の直線に順中心はないので、
第二曲線は位置0のねじれ三次曲線である。これは $z$ の後続点 $z^+$ を頂点とする
二次錐に入る。錐上のねじれ三次曲線は必ず頂点を通る。
避けるなら頂点からの射影の三座標は基点のない三次式であり、滑らかな底二次曲線を
経由する射の引戻し次数が奇数となるためである。
ゆえに $z^+$ のある前点が $a$ の後続直線上にある。位置0に逆中心がなく
$z\to z^+$ なのでその前点は $z$ 自身である。しかし最初の後続直線の全点の前点は一意であり、
逆中心 $z$ はその上に乗れない。

これで順中心は位置0の $p$ と位置1の $q$ のみである。後者の一意性は
全て $L$ 上にあり $f_1=1$ であることによる。逆中心は列の対応の全単射により二つで、
$z$ と、上の定理を逆に使って得る位置2の逆中心 $b$ で尽きる。
$p$ の列では $f_2=f_3=0$、従って $d_3=3$, $d_4=4+b_2\ge4$。
非直線の次も動くので第五曲線がある。第五曲線は位置1に戻るが、最初の $L$ を再訪できない。
次が定点なら位置2の逆中心 $b$ であり、逆順の第二次数が $d_4\ge4$ になって矛盾する。
次が動けば $L$ の外から来るので第六曲線は $N$ で、その次が $z$ となる。
この場合には二項中心 $z$ から逆順の第三次数が $d_4\ge4$ となるが、上の定理では3以下である。
よって二項中心は存在せず、動く点の主張は定理\ref{thm:branch-binary-center-pairing}から従う。
\end{proof}
